% !TeX root = ../report.tex
\newcolumntype{Y}{>{\raggedright\arraybackslash}X}  % wrapping, ragged-right
\newcolumntype{Q}{>{\centering\arraybackslash}X}    % wrapping, centred

\begin{appendices}

% Table numbering A1, A2, ... (restarts inside the appendix)
\setcounter{table}{0}
\renewcommand{\thetable}{A\arabic{table}}
\renewcommand{\theHtable}{A\arabic{table}}  % stops hyperref duplicate-anchor warnings

\section{Global rule extraction experiments}
\label{app:global}

Global rules were extracted once for all 379 drugs together (the ``global'' runs). One search per network and per extractor produced rules over the pooled training data. These runs compared the Adam baseline network with the SAM+WD network. They used the same chromosome (572 features, Low/Intermediate/High) and the same fitness definition as the per-drug experiment, but a different data split and a different coverage floor (Appendix~\ref{app:global-diff}). The extractor settings were otherwise those of Sections~\ref{sec:shared}, \ref{sec:pdere-ga} and~\ref{sec:pdere-nsga}.

\subsection{Setup}
\label{app:global-setup}

\begin{table}[H]
\centering
\small
\caption{Settings of the global runs.}
\label{tab:A1}
\begin{tabularx}{\textwidth}{@{}lY@{}}
\toprule
\textbf{Item} & \textbf{Setting} \\
\midrule
Networks & Baseline: Adam, train--validation gap 0.3369. SAM: $\rho = 0.5$, weight decay $10^{-3}$, gap 0.0485 \\
Data & 245,845 training rows (the 379-drug set); each network's own train/validation/test split \\
Runs & 30 per network and extractor; seeds 43 to 72 (base seed 42 plus run number); none excluded; 50 generations per run \\
Coverage floor & A rule must cover at least 2\% of the training rows \\
Rule pools & 7 (baseline) and 6 (SAM) rules from PDERE-GA; 3 and 3 from PDERE-NSGA; 19 in total \\
\bottomrule
\end{tabularx}
\end{table}

\subsection{Results}
\label{app:global-results}

\begin{table}[H]
\centering
\small
\caption{Fidelity of the extracted rules.}
\label{tab:A3}
\begin{tabularx}{\textwidth}{@{}p{0.27\textwidth}QQQQ@{}}
\toprule
\textbf{Measure} & \textbf{PDERE-GA baseline} & \textbf{PDERE-GA SAM} & \textbf{PDERE-NSGA baseline} & \textbf{PDERE-NSGA SAM} \\
\midrule
Best-rule training fidelity & 0.4800 & 0.4914 & 0.4617 & 0.4553 \\
Median best fidelity per run (30 runs) & 0.4691 & 0.4704 & 0.4066 & 0.4067 \\
Baseline against SAM (Mann--Whitney $p$) & \multicolumn{2}{c}{0.234} & \multicolumn{2}{c}{0.663} \\
\bottomrule
\end{tabularx}
\end{table}

\begin{table}[H]
\centering
\small
\caption{Rule-level tests.}
\label{tab:A4}
\begin{tabularx}{\textwidth}{@{}YYccc@{}}
\toprule
\textbf{Comparison} & \textbf{Measure} & \textbf{Group 1} & \textbf{Group 2} & \textbf{$p$} \\
\midrule
Baseline against SAM & Test fidelity (median) & 0.4706 & 0.4794 & 0.488 \\
Baseline against SAM ($n = 10, 9$) & PROMETHEE~II net flow (median) & 0.0418 & 0.0892 & 0.838 \\
\bottomrule
\end{tabularx}
\end{table}

PROMETHEE~II used the weights fidelity 0.5, length 0.2 and stability 0.3. (In the per-drug experiment stability was analysed separately, so only fidelity and length were ranked, with the weights rescaled to 0.714 and 0.286; Section~\ref{sec:rule-eval}.) The top-ranked rule was \texttt{ga\_sam\_0002} (net flow 0.2776) and the lowest was \texttt{nsga\_sam\_0003} ($-0.4758$). In the gene set enrichment analysis, no adjusted $p$-value was below 0.05.

In short, the global runs found no detectable difference between the baseline and SAM networks.

\subsection{How these runs differ from the per-drug experiment}
\label{app:global-diff}

% No caption: the Google Doc version has none for this table
\begin{table}[H]
\centering
\small
\begin{tabularx}{\textwidth}{@{}lYY@{}}
\toprule
\textbf{Aspect} & \textbf{Global runs} & \textbf{Per-drug experiment} \\
\midrule
Splits & Row-level train/validation/test split of each network & Five folds by cell line \\
PDERE-GA accuracy term & The network's validation split & An inner 20\% of the training folds \\
Coverage floor & 2\% of the training rows & 25 cell lines per drug \\
Scope & One search over all 379 drugs & One search per drug \\
Networks & Baseline and SAM & Adam, Adam+WD and SAM+WD \\
\bottomrule
\end{tabularx}
\end{table}

Fidelity is therefore computed on different rows, and the two sets of numbers should not be compared directly.

\subsection{Limitations}
\label{app:global-limits}

\begin{itemize}
  \item The PROMETHEE and Mann--Whitney tests on rules treat rules pooled across runs as independent samples, but the same rule is often found by several runs. The per-run tests (30 against 30) do not have this problem.
  \item PDERE-NSGA returned a front of 40 rules per run but only 2 to 6 distinct ones, so its rule pools are small.
  \item The PDERE-GA fitness weights and the PROMETHEE weights were set by the team and were not tuned.
\end{itemize}

The rule pools and settings of the global runs are in the project repository (Appendix~\ref{app:data-code}).

\section{Data and results availability}
\label{app:data}

Everything needed to inspect or rerun the study is listed here.

\subsection{Full per-drug extraction results}
\label{app:data-results}

The complete per-drug experiment has 11,370 cells (379 drugs $\times$ 3 networks $\times$ 2 extractors $\times$ 5 folds): 5,685 PDERE-GA cells and 5,685 PDERE-NSGA cells. Each cell holds the extracted rules, their fidelity and accuracy on the held-out cell lines, and the settings used. The results are stored in three archives, each with a \texttt{run\_provenance.json} file that records the command and settings of the run.

% No caption: the Google Doc version has none for this table
\begin{table}[H]
\centering
\small
\begin{tabularx}{\textwidth}{@{}llcY@{}}
\toprule
\textbf{Archive} & \textbf{Extractor} & \textbf{Drugs} & \textbf{Link} \\
\midrule
\texttt{ga\_top\_rev} & PDERE-GA & 189 & \href{https://drive.google.com/file/d/1vZ9ZFul-KmKWqnsSMWdWHV4CVNh_n3Ss/view?usp=drive_link}{ga\_top\_rev} \\
\texttt{ga\_bottom\_rev} & PDERE-GA & 190 & \href{https://drive.google.com/file/d/1lAGty4Nq5CPfCwCa-RQTd6lFtC17cId4/view?usp=drive_link}{ga\_bottom\_rev} \\
\texttt{nsga\_final} & PDERE-NSGA & 379 & \href{https://drive.google.com/file/d/1L1wWk8dy6c1hnCI6QewStOohlpbfy7gI/view?usp=drive_link}{nsga\_final} \\
\bottomrule
\end{tabularx}
\end{table}

All three are in the Drive folder \href{https://drive.google.com/drive/folders/1YIIrkWfg-eUVOTwJKQsFeaIrR9JQLwUV?usp=sharing}{APPENDIX}. To analyse them, unzip the archives into one folder that holds a \texttt{cells/} directory and follow the run instructions in the repository README (\texttt{verify\_cells.py}, then \texttt{analyse.py}).

\subsection{Chromosome feature files}
\label{app:data-features}

Each PDERE chromosome has 572 positions, one per input feature of the drug-response network: 38 KEGG pathway activity features (positions 1 to 38) followed by 534 gene mutation features (positions 39 to 572). A position holds $-1$ (feature not used) or a level. Pathway scores are z-scored across the 878 cell lines and divided into Low ($z < -1$), Intermediate and High ($z > 1$). A mutation feature is High if the gene is mutated and Low if it is not. A pathway feature describes a whole pathway, not one gene.

The full list of features, in chromosome order, is in the files below:

\begin{itemize}
  \item \texttt{chromosome\_features\_mutation\_genes.csv} (534 genes): \href{https://drive.google.com/file/d/1ba-mnUhJxmIIog-WZMatWdbVAE-Zwoh8/view?usp=drive_link}{Link}
  \item \texttt{chromosome\_features\_pathways.csv} (38 pathways): \href{https://drive.google.com/file/d/1ucFABj0zOMYXg7lMVOGzzt2E3JNwWa-T/view?usp=drive_link}{Link}
  \item \texttt{README\_chromosome\_features.txt}: \href{https://drive.google.com/file/d/1DHbwAa4ZrXldvbWXHumVmMFg0RpIrbvL/view?usp=drive_link}{Link}
\end{itemize}

\subsection{Code}
\label{app:data-code}

All code, configuration files and the run instructions are in the project repository: \href{https://github.com/HawaIbr01/COS700-XAI-In-Cancer-Medicine-Treatment}{GitHub}.

\end{appendices}